\documentclass[11pt]{article}
\usepackage[T1]{fontenc}
\usepackage{textcomp}
\usepackage[margin=0.75in]{geometry}
\usepackage{amsmath,amssymb,amsthm}
\usepackage{array,booktabs}
\usepackage{hyperref}
\setlength{\parindent}{0pt}
\newtheorem{theorem}{Theorem}
\theoremstyle{definition}
\newtheorem{definition}{Definition}
\title{Flow Proof Artifact: prop-22-cyclic-quadrilateral-opposite-angles}
\author{Generated by \texttt{flow doc proof}}
\date{Flow Proof Book}
\begin{document}
\maketitle
The opposite angles of quadrilaterals in circles equal two right angles.\\[0.5em]
\textbf{Source.} euclid — Elements, Book III, Proposition 22\\[0.5em]
\bigskip
\noindent{\large \textbf{Derived fact 1.} \textit{The opposite angles of quadrilaterals in circles equal two right angles} \hfill the Euclidean plane \cdot Euclid Book III \cdot Proposition 22: opposite angles of a cyclic quadrilateral sum to two right angles}\par
\smallskip\noindent\textit{Coordinate: the Euclidean plane \cdot Euclid Book III \cdot Proposition 22: opposite angles of a cyclic quadrilateral sum to two right angles}\par
\smallskip\noindent\textit{Source: euclid — Elements, Book III, Proposition 22}\par
\smallskip\noindent\textit{Built on: proposition 21: angles in the same segment are equal, for Euclid Book III on the Euclidean plane, proposition 13: adjacent angles on a straight line sum to two right angles, for Book I of the Elements in on the Euclidean plane}\par
\medskip
\noindent\textbf{Goal.} The opposite angles of quadrilaterals in circles equal two right angles\par
\noindent$\displaystyle \text{opposite angles sum two right}$\par
\medskip
\noindent
\begin{tabular}{@{} >{\raggedright\arraybackslash}p{0.06\textwidth} >{\raggedright\arraybackslash}p{0.47\textwidth} >{\raggedleft\arraybackslash}p{0.41\textwidth} @{}}
\textbf{\#} & \textbf{Proof} & \textbf{Mathematics} \\
\midrule
\textbf{1.} & We prove that proposition 22: opposite angles of a cyclic quadrilateral sum to two right angles for Euclid Book III the Euclidean plane. & \\[0.45em]
\textbf{2.} & Let ABCD = cyclic\_quadrilateral. & \\[0.45em]
\textbf{3.} & Let angle\_A = angle\_DAB. & \\[0.45em]
\textbf{4.} & Let angle\_C = angle\_BCD. & \\[0.45em]
\textbf{5.} & From step 2, step 3, and step 4, we can deduce that angle A plus angle C equals two right angles. & $\displaystyle \angle A + \angle C = 180^\circ$ \\[0.45em]
\textbf{6.} & We invoke the derived fact governing Euclid Book III the Euclidean plane: proposition 21: angles in the same segment are equal, for Euclid Book III on the Euclidean plane. & \\[0.45em]
\textbf{7.} & We invoke the axiom governing Book I of the Elements in the Euclidean plane: proposition 13: adjacent angles on a straight line sum to two right angles, for Book I of the Elements in on the Euclidean plane. & \\[0.45em]
\textbf{8.} & From step 6 and step 7, this implies opposite angles sum two right. Hence proven. & $\displaystyle \text{opposite angles sum two right}$ \\[0.45em]
\bottomrule
\end{tabular}
\label{thm:Geometry-Euclid Book III-proposition_22_opposite_angles_of_a_cyclic_quadrilateral_sum_to_two_right_angles}
\par\medskip\hrule
\end{document}
